% Capitulo: Self-service e a era dos dashboards
\chapter{Self-service e a era dos dashboards}
\label{cap:self-service}

O data warehouse resolveu a consistência, mas manteve a distância. Para ter um relatório
novo, o analista de negócio ainda abria um chamado, esperava a fila da TI e recebia, semanas
depois, algo parecido com o que tinha pedido. Nos anos 2000 e 2010, uma nova geração de
ferramentas encurtou essa distância de forma radical: o próprio analista passou a
conectar os dados, arrastar campos e montar o seu painel. Este capítulo trata dessa
virada, dos problemas que ela criou e do que vale aprender sobre visualização
independentemente da ferramenta.

\section{O BI sai da TI}
\label{sec:bi-sai-da-ti}

Uma das origens dessa mudança é um projeto acadêmico. Na Universidade Stanford, Chris
Stolte, Diane Tang e Pat Hanrahan desenvolveram o Polaris, um sistema que traduzia
gestos visuais (arrastar uma dimensão para as linhas, uma medida para as colunas) em
consultas ao banco de dados e mostrava o resultado como gráfico
\parencite{stolte2002}. A ideia central era que explorar dados devia ser um processo
visual e iterativo, e não a redação de uma consulta seguida da espera pelo resultado. O
projeto deu origem a uma das ferramentas de visualização mais conhecidas do mercado.

Na mesma época, outras ferramentas apostaram em carregar os dados na memória do
computador e permitir que o usuário navegasse livremente por eles, clicando num valor e
vendo todo o resto se reorganizar em torno dele. E, já na metade dos anos 2010, os
grandes fornecedores de software corporativo entraram com ferramentas próprias, muitas
vezes integradas à planilha que o analista já usava todo dia, a preços que tornaram a
licença acessível para qualquer área.

O efeito somado foi uma mudança de quem fazia BI. O analista de negócio, que antes
consumia relatórios, passou a produzi-los. A TI deixou de ser a fábrica de relatórios e
passou, em tese, a cuidar da infraestrutura e dos dados de base. O tempo entre a
pergunta e o primeiro gráfico caiu de semanas para horas.

Foi uma conquista real. Muita gente que hoje trabalha com dados entrou na área por esse
caminho: começou montando painéis para a própria equipe, descobriu que gostava e foi
aprendendo SQL, modelagem e estatística no meio do caminho. A trilha tradicional de
aprendizado que o capítulo~\ref{cap:trilha-tradicional} descreve tem muito desse
período.

\section{Mil dashboards e nenhuma verdade}
\label{sec:mil-dashboards}

A liberdade trouxe de volta o problema que o data warehouse tinha tentado resolver. Cada
área passou a conectar a ferramenta onde achava melhor: no armazém, numa extração, numa
planilha que alguém mantinha à mão. Cada uma aplicava os próprios filtros e as próprias
regras. Não demorou para aparecerem três versões do faturamento do mês, todas
\enquote{oficiais}, todas diferentes.

O sintoma clássico é a reunião em que a primeira meia hora se gasta discutindo de onde
veio cada número, e a decisão fica para a próxima semana. O problema raramente é a
ferramenta. É que a definição da métrica (o que entra no faturamento, se devolução
desconta, em que data a venda conta) foi feita várias vezes, por pessoas diferentes, sem
que ninguém soubesse.

Há também o problema do volume. Como criar um painel ficou fácil, criar mais um painel
passou a ser a resposta para qualquer pergunta. Empresas acumularam centenas de painéis,
muitos abertos uma vez e esquecidos, outros duplicados com pequenas variações, todos
consumindo atualização, manutenção e, cada vez mais, dinheiro. É a crítica de Ackoff
(\veja{sec:sad}) com ferramenta moderna. O capítulo~\ref{cap:frequencia} volta a esse
ponto com uma proposta concreta: medir o uso dos próprios produtos de dados e desligar o
que não é usado.

\rascunho{
  Relato: uma situação de números divergentes entre áreas, ou de painéis duplicados, e
  como foi resolvida. Sem nomes de empresa nem valores exatos.
}

\section{Governança sem burocracia}
\label{sec:governanca}

A primeira reação de muitas empresas ao caos foi recentralizar: tirar a liberdade das
áreas e voltar a fila da TI. Funciona por algum tempo, até as áreas voltarem às
planilhas, porque a fila continua lenta e a pergunta continua urgente. Governança que só
proíbe empurra o trabalho para fora do alcance dela.

O que costuma funcionar é separar o que precisa ser único do que pode ser livre. O que
precisa ser único são os dados de base e as definições das métricas que a empresa usa
para decidir: faturamento, margem, cliente ativo. Esses ficam num lugar só, com dono e
com regra escrita, e as ferramentas de self-service consomem dali. O que pode ser livre é
a exploração em cima dessa base: cada área monta as análises e os painéis de que
precisa, desde que parta dos dados certificados.

Na prática, isso se traduz em conjuntos de dados certificados, que as ferramentas
marcam como oficiais, e numa camada semântica, onde cada métrica é definida uma vez e
reaproveitada por todos os painéis (\veja{sec:camada-semantica}). Aqui o livro só
apresenta o problema. A governança de dados como disciplina, e a forma como a IA muda o
papel de quem cuida dela, ficam para o capítulo~\ref{cap:governanca}.

\section{Visualização que não depende de ferramenta}
\label{sec:visualizacao}

As ferramentas de self-service tornaram fácil fazer gráfico. Não tornaram fácil fazer
gráfico bom. Algumas ideias, todas anteriores às ferramentas atuais, continuam sendo o
melhor ponto de partida.

Edward Tufte propôs que a maior parte da tinta de um gráfico deveria representar dados,
e não decoração: bordas, sombras, fundos e efeitos tridimensionais disputam a atenção
com a informação \parencite{tufte1983}. O princípio vale para a tela. Um painel limpo,
com poucos elementos e muito espaço, comunica mais do que um painel cheio de enfeites.

Stephen Few levou essas ideias para os dashboards e propôs uma definição exigente: um
dashboard mostra, numa única tela, a informação mais importante para atingir um ou mais
objetivos, de forma que possa ser acompanhada de relance \parencite{few2006}. Se o painel
precisa de rolagem, de dez abas e de um manual, ele pode ser útil, mas não é um
dashboard nesse sentido. Few também insistia em comparar sempre: um número sozinho não
diz se está bom ou ruim; o mesmo número ao lado da meta, do mês anterior ou do mesmo mês
do ano passado diz.

Cole Nussbaumer Knaflic acrescentou a perspectiva da comunicação: antes de escolher o
gráfico, é preciso saber quem vai ler, o que essa pessoa precisa decidir e qual é a
mensagem principal \parencite{knaflic2015}. Na maior parte das vezes, uma barra simples,
com o ponto importante destacado e um título que diz a conclusão, funciona melhor do que
qualquer gráfico sofisticado.

\begin{dica}
Antes de publicar um painel, faça três perguntas: quem vai abrir, com que frequência, e
que decisão essa pessoa toma com o que está vendo. Se não houver resposta para a terceira,
o painel provavelmente não precisa existir.
\end{dica}

Essas três referências têm em comum o fato de que nenhuma depende de ferramenta. É
por isso que valem mais do que qualquer curso de um produto específico: o produto muda a
cada poucos anos, o olho de quem lê o gráfico não.

O próximo capítulo volta para os bastidores. Enquanto as ferramentas de visualização
democratizavam a análise, o volume de dados crescia, o armazenamento ficava barato e o
data warehouse começava a se mudar para a nuvem.

\begin{leituras}
  \leitura{few2006}{O livro que definiu o que é um bom dashboard. Os exemplos de painéis
    ruins, analisados um a um, ensinam tanto quanto as recomendações.}
  \leitura{knaflic2015}{Para quem precisa apresentar dados a gestores: como escolher o
    gráfico, tirar o excesso e conduzir a atenção para a mensagem.}
  \leitura{tufte1983}{O clássico da visualização de dados. Mais conceitual que prático,
    mas forma o olhar para sempre.}
  \leitura{stolte2002}{Para ver, na fonte, a ideia de exploração visual que está por trás
    de boa parte das ferramentas de self-service.}
\end{leituras}
